\section*{Boolean payment of the complete singleton nonroot KL layer}

All logarithms are base two.  Let $f:\{-1,1\}^n\to\{-1,1\}$ be
Boolean, put $p(S)=\widehat f(S)^2$, and
$\pi(A)=\sum_{S\supseteq A}p(S)$.  For a coordinate $i$ and
$A\subseteq[n]\setminus\{i\}$ write
\[
 a_{i,A}=\widehat f(A),\quad b_{i,A}=\widehat f(A\cup\{i\}),\quad
 m_{i,A}=a_{i,A}^2+b_{i,A}^2,
\]
\[
 M_{i,A}=\pi(A),\qquad
 C_{i,A}=\sum_{T\supseteq A,\ i\notin T}
              \widehat f(T)\widehat f(T\cup\{i\}),
\]
and
\[
 \alpha_{i,A}={2a_{i,A}b_{i,A}\over m_{i,A}},\qquad
 \beta_{i,A}={2C_{i,A}\over M_{i,A}}.
\]
The row contribution in the exact signed restriction identity is
\[
 K_{i,A}=m_{i,A}D_2\left({1+\alpha_{i,A}\over2}\middle\|
                           {1+\beta_{i,A}\over2}\right).       \tag{1}
\]

For $k\ge0$ put $W_k=\sum_{|S|=k}p(S)$ and retain the sharp root reserve
\[
 M_1(f)=\sum_i\min\{p(\varnothing),p(\{i\})\}.          \tag{2}
\]
For ordered distinct pairs $(i,j)$ define
\[
 m_{ij}=p(\{j\})+p(\{i,j\}),\qquad P_j=\pi(\{j\}),
\]
and the singleton-row self-information
\[
 V_1(f)=\sum_{i\ne j}m_{ij}\log_2{P_j\over m_{ij}},     \tag{3}
\]
with zero-mass terms omitted.

\paragraph{Theorem (complete $|A|=1$ payment).}
For every genuine Boolean $f$,
\[
 \boxed{
 K^{(1)}(f):=\sum_i\sum_{|A|=1}K_{i,A}
 \le {1\over2}V_1(f)+{2\over\ln2}W_2(f).}             \tag{4}
\]
The constants are independent of dimension, bias, arity, overlap and
coordinate reuse.

\paragraph{Proof.}
Fix ordered $i\ne j$, put $A=\{j\}$, and abbreviate $a=a_{i,A}$,
$b=b_{i,A}$, $m=a^2+b^2$, $M=P_j$.  Let
\[
 x_\pm=(a\pm b)^2,
 \qquad X_\pm=M\pm2C_{i,A},
 \qquad p_\pm={x_\pm\over2m},\quad q_\pm={X_\pm\over2M}.
\]
The two functions $f_{i,+}$ and $f_{i,-}$ are again Boolean.  The
singleton case of the Boolean granularity lemma, applied to coordinate
$j$ in each child, gives
\[
 \sqrt{x_\pm}=|\widehat {f_{i,\pm}}(\{j\})|
 \le\operatorname{Inf}_j(f_{i,\pm})=X_\pm.             \tag{5}
\]
Therefore, whenever $p_\pm>0$,
\[
 {p_\pm\over q_\pm}
 ={x_\pm M\over mX_\pm}
 \le M\sqrt{{2p_\pm\over m}}.
\]
Averaging the logarithm under $p$ yields
\[
 D_2(p\|q)
 \le {1\over2}\log_2{M^2\over m}+{1\over2}
       -{1\over2}h_2(p_+)
 \le {1\over2}\log_2{M\over m}
       +{1\over2}\bigl(1-h_2(p_+)\bigr),              \tag{6}
\]
because $M\le1$.  The sharp fair-posterior scalar estimate gives
\[
 m\bigl(1-h_2(p_+)\bigr)
 \le {2\over\ln2}\min\{a^2,b^2\}
 \le {2\over\ln2}p(\{i,j\}).                          \tag{7}
\]
Multiplying (6) by $m$ and applying (7),
\[
 K_{i,\{j\}}
 \le {1\over2}m_{ij}\log_2{P_j\over m_{ij}}
       +{1\over\ln2}p(\{i,j\}).                       \tag{8}
\]
Summation over ordered pairs counts each two-set twice and proves (4).
No unsigned replacement of $C_{i,A}$ was made; it remains inside $q_\pm$
until the genuine Boolean child estimate (5) is used.

\paragraph{Exact GU consequence.}
Let
\[
 D_{\ge2}(f)=\sum_i(B_i-\Gamma_i)
 +\sum_i\sum_{|A|\ge2}K_{i,A}.                         \tag{9}
\]
The exact local identity and the sharp root theorem give
\[
 G=D_{\ge2}+K^{(0)}+K^{(1)}
 \le D_{\ge2}+{1\over2}V_1
       +{2\over\ln2}(M_1+W_2).                        \tag{10}
\]
Since $M_1\le W_1$ and
$M_1+W_2\le W_1+W_2\le I$, one obtains the uniform reduction
\[
 \boxed{G(f)\le D_{\ge2}(f)+{1\over2}V_1(f)
                    +{2\over\ln2}I(f).}               \tag{11}
\]

Freeze $R(f)=G(f)-4I(f)-\tfrac34U(f)$.  If the frozen GU pair fails,
choose a counterexample of minimum dimension and a tensor-prime
positive-residual factor.  Every proper physical child has nonpositive
residual, and hence every relevant coordinate satisfies
\[
 \Delta_iG-\tfrac34\Delta_iU
 \ge4\operatorname{Inf}_i(f)+R(f).                    \tag{12}
\]
No optimizer over unbounded dimension is assumed.  In view of (11), it is
now sufficient, only on tensor-prime Boolean functions with $R>0$, to prove
\[
 \boxed{
 D_{\ge2}(f)+{1\over2}V_1(f)
 \le\left(4-{2\over\ln2}\right)I(f)+{3\over4}U(f).}   \tag{13}
\]
Then $G\le4I+3U/4$ universally, so restriction induction gives
$U\le16I$ and the source reduction gives $H\le34I$.

Compared with the preceding endpoint, (13) removes every nonempty
singleton KL row from the unpaid signed aggregate.  The remaining KL sum
starts at $|A|=2$.  The explicit term $V_1/2$ is retained; no claim that it
is already bounded by $I$ or $U$ is made.
